\documentclass[11pt,a4paper]{article}
\usepackage[margin=0.85in,top=1in,bottom=1in]{geometry}
\usepackage{amsmath,amssymb,amsfonts}
\usepackage{graphicx}
\usepackage{float}
\usepackage{enumitem}
\usepackage{microtype}
\usepackage{caption}
\usepackage[hidelinks]{hyperref}
\usepackage[most,breakable]{tcolorbox}
\tcbuselibrary{breakable,skins}
\usepackage[utf8]{inputenc}
\usepackage[T1]{fontenc}
\usepackage{lmodern}
\captionsetup{font=small,labelfont=bf,skip=4pt}
\setlist{leftmargin=1.5em,topsep=2pt}

%% Breakable card — prevents content cutoff across pages
\newtcolorbox{solcard}[3]{
  breakable, enhanced,
  colback=white, colframe=gray!50, arc=2pt,
  title={\textbf{#1}\hfill\textsf{\small #3\,pts}},
  fonttitle=\normalsize\sffamily\bfseries,
  before upper={\small\textit{#2}\par\smallskip\hrule height 0.4pt\smallskip},
  top=4pt, bottom=6pt, left=7pt, right=7pt,
  parbox=false,
}

\title{\textbf{Microfluidics --- Solution}}
\author{\normalsize X23 (2023) — English translation}
\date{}
\begin{document}\maketitle\vspace{-0.5em}
\begin{solcard}{A1}{Let there be a steady laminar flow of water in the capillary. Find the dependence of the liquid flow velocity $v(z)$ in the capillary on the coordinate $z$ for steady laminar flow of water in the capillary. Axis $z$ is directed along side $b$, the origin of coordinates is on one of the walls of the capillary. Express your answer in terms of $\Delta P$, capillary parameters and liquid viscosity $\eta$.}{1.00}

Let us consider a thin layer of liquid enclosed between planes with coordinates $z_1=z$ and $z_2=z+dz$. Viscous friction forces acting in the corresponding planes: $$ f_1 = -\eta \ aL \frac{dv_x}{dz} (z) \\ f_2 = \eta \ aL \frac{dv_x}{dz} (z+dz) $$ Then Newton's second law assume form: $$ \Delta Pa \ dz + f_1+f_2 = 0 \;(\text{1}) $$ Hence equation on velocity: $$ \frac{d^2v_x}{dz^2} = -\frac{\Delta P}{\eta L} \;(\text{2}) $$ Note, that completely not necessarily be take to the limit thickness layer to zero (Thus, equation (1) correct and for final $dz$). More that, record (for example $z_1 =0$ or $z_1 = b/2$), we obtain integrate analog equation (2). General solution equation (2) have form: $$ v_x = -\frac{\Delta P}{\eta L} \frac{z^2}{2} + C_1 z + C_2 $$ Constant integration be determined boundary condition, in given case – $v(0)=0$ and  $v(b)=0$. From here we get the correct answer

\par\smallskip\noindent\textbf{Answer:} $$ v(z) = \frac{\Delta P}{2\eta L} \left(bz -z^2 \right) $$
\end{solcard}

\begin{solcard}{A2}{At some point, a reagent is uniformly injected into the initial section of the capillary $AA'$. In the figure, thick horizontal lines indicate the walls of the capillary. On your answer sheet, draw where the reactant will be located via $\tau=5~\mathrm{s}$.  For the drawing, select an appropriate scale along the horizontal axis, which is directed along the fluid flow. Capillary parameters: $L = 8~\text{mm}$, $b = 0.2 ~\text{mm}$, pressure difference $\Delta P = 1.5 ~\text{Pa}$, $\eta = 8.9 \cdot 10^{-4}~\text{Pa} \cdot \text{s} $. Neglect the diffusion of the reagent.}{1.00}

Let's plot the dependence $x(z) = \tau v(z)$ on the graph

\begin{figure}[H]
  \centering
  \includegraphics[width=0.72\textwidth]{T1_S_EN_assets/fig1.jpg}
\end{figure}
\end{solcard}

\begin{solcard}{B1}{Find the steady-state velocity of ions $v_z(z)$ at a distance $z$ from the capillary boundary. Assume that the movement of the ion during the establishment of the velocity is small. Express your answer in terms of the charge of the ion $q(q = \pm e)$, $\alpha$ and the potential $\varphi(z)$ (and its derivatives).}{0.50}

In steady state, Newton's second law in projection onto the $z$ axis looks like this: $$ -\alpha v_z + q E_z = 0 $$ Express electric field through gradient potential $E_z = - \tfrac{d \varphi}{dz}$, we get the answer

\par\smallskip\noindent\textbf{Answer:} $$ v_z = - \frac{q}{\alpha} \frac{d \varphi}{dz} \;(\text{3}) $$
\end{solcard}

\begin{solcard}{B2}{Express the concentrations of cations ($n_+$) and anions ($n_-$) near the capillary wall in terms of the potential at this point and also in terms of $n_0, e, \alpha, k, T$.}{1.50}

The flow of particles associated with the presence of an electric field is balanced by the flow of ions associated with diffusion. Thus, the balance condition will look like this: $$ n v_z -D \frac{dn}{dz} = 0 $$ Substitute here equation (3), obtain $$ -\frac{q}{\alpha D} \frac{d \varphi}{dz} = \frac{1}{n} \frac{dn}{dz} $$ Integrate both part equation by $z$ and substitute $D=\frac{kT}{\alpha}$, find $$ -\frac{q}{kT} (\varphi - 0) = \ln(\frac{n}{n_0}) $$ As a result,


The result of this paragraph is a special case of the Boltzmann distribution. Of course, this coincidence is not accidental.

\par\smallskip\noindent\textbf{Answer:} $$ n = n_0 \exp \left( -\frac{q \varphi}{kT} \right) $$
\end{solcard}

\begin{solcard}{B3}{Express the volume density of the uncompensated charge $\rho(z)$ near the capillary wall in terms of potential $\varphi(z)$, ion concentration $n_0$ and physical constants $e, k, T$. Assuming that the potential applied to the wall is small ($e\varphi_0 \ll kT$), you will obtain an approximate value of $\rho(z)$.}{0.50}

$$ \rho = e(n_+ - n_-) = en_0 \left[ \exp \left( -\frac{e \varphi}{kT} \right)  - \exp \left( \frac{e \varphi}{kT} \right) \right] = - 2 e n_0 \sinh{\left( \frac{e \varphi}{kT} \right)} \;(\text{4}) $$ Expand expression by parameter $\frac{e \varphi}{kT}$, we find

\par\smallskip\noindent\textbf{Answer:} $$ \rho \approx - 2 \frac{e^2 n_0 \varphi}{kT} \;(\text{5}) $$
\end{solcard}

\begin{solcard}{B4}{Find the dependence of the potential $\varphi(z)$ on the coordinate near the capillary wall in the diffuse layer. Remember that the potential at the boundary is $\varphi_\delta$. Express your answer in terms of $n_0$, $\varphi_\delta$, $T$, physical constants and electrolyte parameters.}{1.00}

Gauss's theorem in differential form in the one-dimensional case has the form: $$ \rho = - \varepsilon \varepsilon_0 \frac{d^2 \varphi}{dz^2} $$ Use expansion (5) obtain $$ \frac{d^2 \varphi}{ dz^2} = 2 \frac{e^2 n_0}{\varepsilon \varepsilon_0 kT} \varphi \;(\text{6}) $$ This linear differential equation, it general solution in given case look like following manner: $$ \varphi(z) = C_1 \exp \left( z/\lambda \right) + C_2 \exp \left( - z/\lambda \right), $$ where $\lambda$ be determined from the equation (6), and coefficient $C_1$ and $C_2$ – from initial condition $\varphi(0) = \varphi_\delta$ and $\varphi(+\infty) = 0$. Note, that the second boundary condition automatically imply $\lambda \ll b$, Thus, finding second wall on infinity. For accounting for deviation from this condition, strictly say, $C_1 \neq 0$. However, the answer to point B5 allows us to cast aside any doubts about this. As a result


\textit{Note.} Let us note a more interesting point from the point of view of physics. For expansion (5) and characteristic values of $\varphi_\delta \sim 1\:  \mathrm{V}$ to be applicable, temperature $T$ must be $ T \gg \frac{e \varphi_\delta}{k} \sim 10^4\: \text{K}$! However, the approximation used in the problem is not as bad as it might seem at first glance. Let us verify this by solving the equations without using approximation (5). From combination (4) and Gauss's theorem we find: $$ \frac{d^2 \varphi}{dz^2} = 2 \frac{e n_0}{\varepsilon \varepsilon_0} \sinh{\left( \frac{e \varphi}{kT} \right)} $$ Further, integrate equation by $\varphi$: $$ \int \sinh{\left( \frac{e \varphi}{kT} \right)} d \varphi = \frac{kT}{e} \cosh{\left( \frac{e \varphi}{kT} \right)} + C \\ \int \frac{d^2 \varphi}{dz^2} d\varphi = \int \frac{d \varphi'_z}{dz} \varphi'_z \: dz = \int \varphi'_z d  \varphi'_z = \frac{ (\varphi'_z)^2}{2} + C, $$ where $\varphi'_z \equiv \frac{d \varphi}{dz} = - E_z $. Substitute boundary condition $\varphi'_z (+\infty) = 0$  obtain: $$ \varphi'_z = -2\sqrt{ \frac{e n_0}{\varepsilon \varepsilon_0} \frac{kT}{e} \left( \cosh{\left( \frac{e \varphi}{kT} \right)} -1 \right)} = -2\sqrt{2\frac{e n_0}{\varepsilon \varepsilon_0} \frac{kT}{e}} \sinh{\left( \frac{e \varphi}{2kT} \right)} $$ Further, separating variable and integrate: $$ -2\sqrt{2\frac{e n_0}{\varepsilon \varepsilon_0} \frac{kT}{e}} z = \int  \limits_{\varphi_\delta}^\varphi \frac{d \varphi}{\sinh{\left( \frac{e \varphi}{2kT} \right)}} = \frac{2kT}{e} \ln{\tanh{\frac{e \varphi}{4kT}}} \Bigg|_{\varphi_\delta}^\varphi $$ Hence $$ \varphi = \frac{4kT}{e} \tanh^{-1} \left[ \tanh{ \left( \frac{e \varphi_\delta}{4kT} \right)}{\exp \left(- z/\lambda \right)} \right], $$ where $\lambda$ in accuracy equal its own value in previous consideration. Now let us analyze the asymptotic behavior of the given solution at infinity. For $z \gg \lambda$ exponential turn out small, therefore in leading approximation $\tanh^{-1} (x) \approx x$, and real dependence turn out $$ \varphi \rightarrow  \frac{4kT}{e} \tanh{ \left( \frac{e \varphi_\delta}{4kT} \right)}{\exp \left(- z/\lambda \right)} $$ Such manner, our approximation on lower temperature give correct value of $\lambda$, however the coefficient before the exponent (prefactor) will be inaccurate.

\par\smallskip\noindent\textbf{Answer:} $$ \varphi(z) = \varphi_\delta \: \exp \left( - z/\lambda \right), \;(\text{7}) $$ where $\lambda = \sqrt{\tfrac{\varepsilon \varepsilon_0 kT}{2e^2 n_0}}$.
\end{solcard}

\begin{solcard}{B5}{Find the distance $\lambda$ (formula and numerical value) at which the potential $\varphi(z)$ decreases by a factor of $e$. This length is considered the characteristic size of the Gouy diffuse layer. Liquid viscosity $\eta = 8.9 \cdot 10^{-4}~\text{Pa} \cdot \text{s} $, ion concentration $n = 1.1 \cdot 10^{25}~\text{m}^{-3}$, electrical constant $\varepsilon_0 = 8.85 \cdot 10^{-12} ~\text{F}/ \text{m}$, Boltzmann constant $k = 1.38 \cdot 10^{-23}~\text{J}/\text{K}$, temperature $T = 300~\text{K}$, dielectric constant $\varepsilon = 100$.}{0.50}

The potential decreases by a factor of $e$ at $z = \lambda$, the formula for $\lambda$ was found in paragraph $B4$.

\par\smallskip\noindent\textbf{Answer:} $\lambda = \sqrt{\tfrac{\varepsilon \varepsilon_0 kT}{2e^2 n_0}} \approx 2,55 \mathrm{nm}$
\end{solcard}

\begin{solcard}{C1}{Write down Newton's second law for a small volume of liquid in projection onto the direction of the electric field, taking into account viscous friction and the electric field acting on the uncompensated charge.  Express the answer in terms of the electric field $E$, the ion concentration $n_0$, the potential $\varphi(z)$, the fluid and capillary parameters, and the velocity $v(z)$ and its derivatives.}{0.80}

This point is similar to the solution to point $A1$. The viscous friction forces acting on the layer will be equivalent, and the electric field will act with a force of $$ F_E = \rho E\: aL \: dz $$ Then similarly equation $(1)$: $$ \eta\frac{d^2v}{dz^2} + \rho E = 0 $$ Taking into account expression $(5)$ and $(7)$: $$ \eta \frac{d^2 v}{dz^2} -  \frac{\varepsilon \varepsilon_0 E}{\lambda^2} \varphi_\delta \exp \left( - z/\lambda \right)= 0 $$

\par\smallskip\noindent\textbf{Answer:} $$ \eta \frac{d^2 v}{dz^2} -  \frac{\varepsilon \varepsilon_0 E}{\lambda^2} \varphi_\delta \exp \left( - z/\lambda \right)= 0 $$
\end{solcard}

\begin{solcard}{C2}{Find the dependence of the fluid flow velocity on the coordinate $v(z)$.}{1.20}

From point $C1$ we have the differential equation: $$ \frac{d^2 v}{dz^2} =  \frac{\varepsilon \varepsilon_0 E}{ \eta \lambda^2} \varphi_\delta e^{-z/\lambda } $$ Find it solution, integrate both part by $dz$ twice: $$ v = \frac{\varepsilon \varepsilon_0 E}{ \eta} \varphi_\delta e^{-z/\lambda } + C_1z + C_2 \quad (C_1, C_2 = const) $$ From boundary condition $v(0) = 0$ find: $$ C_2 = -\frac{\varepsilon \varepsilon_0 E}{ \eta} \varphi_\delta $$ With second boundary condition still be the case somewhat complex. Strictly say, $\textit{correct}\:$ boundary condition be $v(b)=0$. However for that, so that consider second wall necessarily introduce change and in distribution potential, add an increasing exponential. But with point of view physics this only add such same picture distribution potential at other wall. For consideration same only one wall arise following issue: using the approximation for $z \sim \lambda$ any boundary condition for $z \sim b$ be indistinguishable from each other, since in fact they still are equivalent infinity and term form $z/b \ll 1$ negligibly are small. The more not less if we we want find distribution potential on scale size tube, the contribution this term becomes atessential. in our same model find $C_2$ one can following manner. Let us turn to formula (8) – it is valid for any $\varphi (z)$ . Thus, boundary condition on wall identical, $v(z)$ and $\varphi (z)$ must be are symmetric about the center of the tube.  From this fact and formula (8) follow $C_1 = 0 $.


\textit{Note.} In fact, it is possible to obtain a more general expression for the speed $v$, showing that it depends only on $\varphi (z)$, while the form of the dependence of the potential on the coordinate itself is unimportant. Let's write the equation: $$ \eta\frac{d^2v}{dz^2} + \rho E = 0 $$ Taking into account theorem at it is written in form: $$ \eta\frac{d^2v}{dz^2} - \varepsilon \varepsilon_0 \frac{d^2 \varphi}{dz^2} E = 0 $$ Similarly two times integrate by $dz$ and obtain: $$ v =\frac{ \varepsilon \varepsilon_0 E \varphi}{\eta} + C_3 z + C_4 \quad (C_3, C_4 = \text{const}). $$ $C_3 = 0$ in force of symmetry tube. For $\varphi = \varphi_\delta$ and $v=0$, from which $C_4 = -\frac{ \varepsilon \varepsilon_0 E \varphi_\delta}{\eta}$. The desired answer: $$ v = \frac{ \varepsilon \varepsilon_0 E}{\eta}(\varphi - \varphi_\delta) \;(\text{8}) $$

\par\smallskip\noindent\textbf{Answer:} $$ v =-\frac{\varepsilon \varepsilon_0 E}{\eta} \varphi_\delta \left( 1- e^{- z/\lambda} \right) $$
\end{solcard}

\begin{solcard}{C3}{Let there be a steady laminar flow of water in the capillary. At some point, a reagent is uniformly injected into the section $AA'$ of the capillary, the pattern is the same as in part $\mathrm{A}$.  On your answer sheet, draw where the reactant will be located through $\tau=5~\mathrm{s}$. Use the capillary parameters from $\mathrm{A2}$ and the numerical data from $\mathrm{B5}$. Potential $\varphi_\delta = 1.8~\text{V}$, voltage drop along the length of the capillary $U = 0.3~\text{V}$. Neglect the diffusion of the reagent outside the diffusion layer. For the drawing, select an appropriate scale along the horizontal axis directed along the direction of fluid flow.}{1.00}

Using the formula for the speed obtained at point $C2$, the displacement of the liquid layer ($\Delta x = v \tau$) is easily calculated. Once you have the values for several $z$, you can plot the graph shown below.


Note that the exponential decays significantly already at $z \sim \lambda$. Considering that $\lambda \ll b$, with $z \sim bv \approx -\frac{\varepsilon \varepsilon_0 E}{\eta} \varphi_\delta = const.$, and in the figure the front will be depicted as a straight line parallel to the $z$ axis.


Photos of the real experiment. On the left is the reagent distribution profile in the case of uniformly distributed pressure, on the right is the case considered in part C.

\begin{figure}[H]
  \centering
  \includegraphics[width=0.72\textwidth]{T1_S_EN_assets/fig2.jpg}
\end{figure}
\begin{figure}[H]
  \centering
  \includegraphics[width=0.72\textwidth]{T1_S_EN_assets/fig3.jpg}
\end{figure}
\end{solcard}

\begin{solcard}{C4}{Based on part $\mathrm{with}$ of this problem, propose a system that allows mixing layers of liquid flowing near the wall of a capillary.}{1.00}
\par\smallskip\noindent\textbf{Answer:} In this part, various systems are possible; we will give a qualitative example of one of them. Let us denote by $Ox$ the axis codirectional with $\vec{E}$ in the figure from part $C$. Let the projection of the field onto this axis have the following form depending on the coordinate $x$: $$ E_x = E_0sin(k x) $$ On base previous part part $C$ we can argue that this will lead to the fact that the layers in some places will be forced to move towards each other. This movement will cause them to mix.
\end{solcard}

\end{document}